\section{Cái trần nói gì và không nói gì về độ trung thực}
\label{sec:ch13-tran-noi-gi}

\index{giới hạn lý thuyết!cái nó không ràng buộc|(}%
Kết luận vừa rút ra trên một chỗ hẹp, chuyện LIME và SHAP, đúng cho cả chương:
cái trần chạm vào độ khớp và không chạm vào câu hỏi của định
nghĩa~\ref{def:ch01-do-trung-thuc}. Dựng lại nó ở mức cả chương là món nợ chương
này mang từ chương~\ref{ch:con-nguoi-vang-mat}. Trước hết là ba điều cái trần
nói được.

Thứ nhất, có một cái sàn, và nó là sàn theo nghĩa mạnh nhất. Với một ngân sách
đọc cố định và một mô hình đủ phức tạp, một lượng sai số nào đó là không tránh
được: không phải khó tránh với công nghệ hiện nay, mà không tránh được, bởi bất
kỳ phương pháp nào đã có hoặc sẽ có. Bốn chương trước đó không có phát biểu nào
thuộc loại này. Chúng nói rằng một dụng cụ đo đã thất bại trên một tập, rằng
lĩnh vực chưa có chuẩn xếp hạng, rằng phép đối chiếu với người hầu như chưa
chạy; mọi câu ấy đều có thể đổi khi có thêm nghiên cứu. Câu ở đây thì không.

Thứ hai, nó nói cái sàn ấy đến từ đâu, và nguồn gốc lại tầm thường tới mức đáng
nói. Định lý~\ref{thm:ch13-khoang-cach-do-phuc-tap} đứng trên tính không nén
được: một chương trình tính $g$ mà $g$ trùng $f$ thì cũng là một chương trình
tính $f$, nên nó không thể ngắn hơn chương trình ngắn nhất tính $f$. Định
lý~\ref{thm:ch13-chan-duoi-lipschitz} và kết quả về hàm ngẫu nhiên đứng trên
một phép đếm: số hàm mô tả được trong $b$ bit không quá $2^{b+1}$, còn số hàm mà
một mô hình có thể là thì nhiều hơn thế rất nhiều. Hai chỗ dựa khác nhau, cùng
một tính chất: không kỹ thuật giải thích nào đổi được chúng, vì không chỗ nào
trong hai chỗ ấy nói gì về kỹ thuật.

Thứ ba, nó cho sự đánh đổi ấy một hình dạng có tốc độ, và đây là điều phải phát
biểu cẩn thận nhất trong ba điều. Câu
\enquote{giải thích mô hình nhiều chiều thì khó} vốn là kinh nghiệm chung của
lĩnh vực; ở đây nó thành một phát biểu có số mũ, mũ theo số chiều và đa thức
theo nghịch đảo ngưỡng sai số. Mục~\ref{sec:ch13-tran-cao-bao-nhieu} cho phát
biểu ấy cả hai chiều: chặn trên nói phép dựng lưới tốn ngần ấy, và phép đếm trên
họ hàm chóp nói không lời giải thích nào tốn ít hơn, cùng số mũ theo số chiều.
Số mũ ấy vì vậy là của bài toán. Và nó cho một cách thoát cụ thể:
số mũ trong chặn trên là chiều nội tại của dữ liệu chứ không phải chiều của
không gian chứa nó, nên
giảm chiều trước khi giải thích không phải mẹo kỹ thuật mà là điều kiện để bài
toán có lời giải trong ngân sách. Đọc theo ngôn ngữ quy định thì cùng ba điều ấy
thành định lý~\ref{thm:ch13-bo-ba-bat-kha}: một bên soạn luật phải chọn bỏ một
trong ba thứ mình muốn, và cái hình~\ref{fig:ch13-mien-kha-thi} vẽ chính là tập
các lựa chọn còn lại sau khi chọn.

Bây giờ tới ba điều nó không nói.

\index{độ trung thực!không nằm trong khung của bài}%
Điều thứ nhất là điều quan trọng nhất, và nó đọc thẳng ra từ định nghĩa. Sai số
mà bài chặn đo trên đầu ra: $\rho(f(x), g(x))$ so hai giá trị mà hai hàm trả về
tại cùng một điểm. Định nghĩa~\ref{def:ch01-do-trung-thuc} hỏi một câu khác, là
lời giải thích có phản ánh đúng quá trình mà mô hình dùng để tới đầu ra hay
không. Hai lời giải thích có thể có sai số đầu ra bằng nhau mà quy công cho
những đặc trưng khác hẳn nhau, và khung của bài không phân biệt được chúng, vì
nó không nhìn vào bên trong $f$ lần nào.

Chỗ ấy còn sắc hơn nếu nhìn vào chính điểm tối ưu của khung. Bài chứng minh
$\kappa_f(0) = K(f)$, nghĩa là lời giải thích duy nhất đạt sai số không là một
chương trình tính đúng hàm $f$. Trong khung này, đó là lời giải thích hoàn hảo.
Với định nghĩa~\ref{def:ch01-do-trung-thuc} thì nó chưa chắc trung thực chút
nào: hai chương trình tính cùng một hàm có thể tính theo hai cách hoàn toàn khác
nhau, và một bảng tra cứu liệt kê mọi cặp đầu vào đầu ra cũng tính đúng hàm ấy
mà không phản ánh bước nào trong quá trình mô hình dùng. Nó lại còn là đúng thứ
mà cả lĩnh vực XAI sinh ra để tránh, vì một chương trình dài bằng mô hình thì
không ai đọc được. Vậy đại lượng mà bài cực tiểu đạt cực trị tại một chỗ vừa vô
dụng với người đọc vừa không có bảo đảm nào về độ trung thực, và điều đó cho
thấy nó không phải đại lượng mà Phần~IV đang tìm.

Còn một chỗ nữa, hẹp hơn và kỹ thuật hơn. Định
nghĩa~\ref{def:ch07-ham-attribution} cho một \tn{hàm attribution}{attribution
function} gán một điểm số cho
mỗi đặc trưng, tức trả về một vector số chứ không phải một đầu ra cùng loại với
đầu ra của mô hình. Một vector như thế không phải một hàm
$\mathcal{X} \to \mathcal{Y}$, nên nó không rơi vào định nghĩa lời giải thích của
bài. Attribution chỉ lọt vào khung ấy qua đường mô hình thay thế: $g$ của LIME
là một hàm đúng kiểu, và mô hình cộng tính của
phương trình~\ref{eq:ch04-mo-hinh-cong-tinh} cũng vậy. \tn{Bản đồ
nhiệt}{heatmap} của Grad-CAM thì không. Cho nên cái trần phủ đúng một nửa Phần~II, và phủ nửa ấy theo tư cách
mô hình thay thế chứ không theo tư cách điểm số.

Điều thứ hai bài không nói là bất cứ điều gì về các họ chỉ số của
chương~\ref{ch:do-do-trung-thuc}. Nó không xác nhận deletion, không bác ROAR,
không nói Sensitivity-n đo đúng hay đo sai. Nó không cần tới chúng để phát biểu
kết quả của mình, và cũng không chạm tới chúng ở đâu. Đây là chỗ mà phép đếm chữ
của mục~\ref{sec:ch13-cau-hoi-khong-can-thi-nghiem} có ích: một người đọc lướt
có thể mang cái trần này về làm bằng chứng rằng lời giải thích nói chung không
đáng tin, và bằng chứng ấy không có trong bài.

\index{giới hạn lý thuyết!và chuỗi kiểm chứng}%
Điều thứ ba là câu trả lời cho
chương~\ref{ch:con-nguoi-vang-mat}, và câu trả lời là không. Chuỗi mà chương ấy
mô tả là chuỗi các dụng cụ đo: mỗi phép kiểm chứng lại đòi một phép kiểm chứng
khác đứng dưới nó, và câu hỏi là chuỗi ấy có đáy hay không. Cái trần của chương
này không phải một dụng cụ đo. Nó không chấm điểm lời giải thích nào, không cần
ground truth, và vì thế không đứng ở đâu trong chuỗi ấy cả. Nó là đáy của một
chuỗi khác, chuỗi hỏi một bản rút gọn có thể tốt tới đâu, và hai chuỗi ấy gặp
nhau ở tên gọi chứ không ở đối tượng.

Nhưng nó thu hẹp được chuỗi kia một chỗ, và chỗ ấy đáng ghi vì
chương~\ref{ch:khoang-trong-nghien-cuu} sẽ cần. Nếu có ngày lĩnh vực dựng được
một dụng cụ đo độ trung thực đã kiểm định, thì dụng cụ ấy không thể có nhiệm vụ
xác nhận rằng một lời giải thích đọc được trùng khớp với một mô hình phức tạp,
bởi định lý~\ref{thm:ch13-khoang-cach-do-phuc-tap} đã nói trước rằng nó không
trùng. Nhiệm vụ còn lại cho một dụng cụ như vậy là mô tả chỗ lệch: lệch ở đâu
trong không gian đầu vào, lệch theo hướng nào, và chỗ lệch ấy có rơi vào vùng mà
người dùng quan tâm hay không. Đó là một yêu cầu cụ thể hơn hẳn cái mà Phần~IV
đã nêu tới giờ, và nó đến từ một định lý chứ không từ một phép đo.

Cuối cùng, khung của bài để lộ một chỗ hổng mà chính bài đặt tên và không theo
tiếp. Điều kiện \tn{không suy biến}{non-degenerate} tồn tại vì nếu thiếu nó thì
một lời giải thích có
thể lệch khỏi mô hình ở rất nhiều điểm mà vẫn giữ sai số dưới ngưỡng, do các chỗ
lệch đều nhỏ hơn ngưỡng. Bài gọi chúng là các bất đồng vô hình và dùng điều kiện
để loại chúng ra khỏi phát biểu. Với một mô hình hồi quy thật thì điều kiện ấy
hỏng, nên các bất đồng vô hình có thật, và chúng là những chỗ mà một lời giải
thích sai lệch mà không chỉ số nào dựa trên sai số đầu ra phát hiện được.

\index{giới hạn lý thuyết!cái nó không ràng buộc|)}%
